% ============================================================
% Chapter 4 — Complete Hardware Design
% Source: PSYCON Engineering Design Document, Vol. 1, Ch. 4
% ============================================================
\section{Complete Hardware Design}
\label{sec:hardware-design}

\subsection{Hardware Design Philosophy}
\label{sec:hardware-design-philosophy}

The hardware platform was designed with the following priorities:
\begin{itemize}
    \item Research-grade sensor selection within a student budget.
    \item Modular architecture for independent testing.
    \item Commercially available components.
    \item Easy replacement and maintenance.
    \item Future migration to a custom PCB.
    \item Low-voltage battery-powered operation.
    \item Standard communication protocols (I\textsuperscript{2}C, I\textsuperscript{2}S, ADC).
\end{itemize}

\subsection{Complete Bill of Hardware}
\label{sec:bill-of-hardware}

\Cref{tab:bill-of-hardware} lists the complete hardware inventory for the PSYCON prototype, including quantity, module assignment, and function of each component.

\begin{table*}[t]
\centering
\caption{Complete bill of hardware for the PSYCON prototype.}
\label{tab:bill-of-hardware}
\small
\begin{tabularx}{\textwidth}{@{}Y c l Y@{}}
\toprule
\textbf{Component} & \textbf{Qty.} & \textbf{Module} & \textbf{Function} \\
\midrule
ESP32 DevKit (30-pin) & 2 & Both & Main controller \\
TP4056 USB-C charger (HW-373 V1.2.1, protected) & 2 & Both & Battery charging and protection \\
500\,mAh 3.7\,V LiPo (HLCY 651735P) & 3 & Both & Power source \\
MAX30102 & 1 & Wrist & Heart rate \& PPG \\
ADS1115 & 1 & Wrist & 16-bit ADC \\
GSR electrodes & 1 pair & Wrist & Skin conductance \\
MPU6050 & 1 & Wrist & Motion sensing \\
MCP9808 & 1 & Wrist & Temperature \\
INMP441 & 1 & Audio & Digital microphone \\
TEMT6000 & 1 & Audio & Ambient light \\
\bottomrule
\end{tabularx}
\end{table*}

\subsection{ESP32 DevKit (30-pin)}
\label{sec:esp32-devkit}

\textbf{Role.} Primary microcontroller responsible for sensor acquisition, data processing, wireless communication, timestamp generation, power management, and firmware execution.

\Cref{tab:esp32-specs} lists the key specifications of the ESP32 DevKit.

\begin{table}[H]
\centering
\caption{Key specifications of the ESP32 DevKit (30-pin).}
\label{tab:esp32-specs}
\small
\begin{tabularx}{\columnwidth}{@{}l Y@{}}
\toprule
\textbf{Parameter} & \textbf{Value} \\
\midrule
MCU & ESP32-WROOM-32 \\
CPU & Dual-core Xtensa LX6 \\
Clock & Up to 240\,MHz \\
SRAM & 520\,KB \\
Flash & Typically 4\,MB \\
Logic voltage & 3.3\,V \\
Wi-Fi & 802.11 b/g/n \\
Bluetooth & Classic + BLE \\
ADC & 12-bit (used sparingly) \\
DAC & 2 channels \\
GPIO & Up to 34 usable \\
\bottomrule
\end{tabularx}
\end{table}

\textbf{Why it was selected.} Integrated Wi-Fi and BLE; large community support; mature software ecosystem; multiple communication interfaces; sufficient processing capability; excellent cost-to-performance ratio.

\textbf{Interfaces used.}
\begin{itemize}
    \item I\textsuperscript{2}C --- MAX30102, MPU6050, ADS1115, MCP9808
    \item I\textsuperscript{2}S --- INMP441
    \item Analog --- optional TEMT6000 input; battery monitoring (future)
\end{itemize}

\begin{engnote}
The ESP32 operates at 3.3\,V logic and its GPIOs are not 5\,V tolerant. The internal ADC is adequate for basic measurements, but the ADS1115 is preferred for precision analog sensing.
\end{engnote}

\subsection{MAX30102}
\label{sec:max30102}

\textbf{Function.} Photoplethysmography (PPG) sensor for cardiovascular monitoring.

\textbf{Measurements.} Heart rate; pulse waveform; SpO\textsubscript{2} capability (not currently a primary feature).

\textbf{Interface.} I\textsuperscript{2}C. \textbf{Default address:} \texttt{0x57}.

\Cref{tab:max30102-elec} lists the electrical characteristics of the MAX30102.

\begin{table}[H]
\centering
\caption{Electrical characteristics of the MAX30102.}
\label{tab:max30102-elec}
\small
\begin{tabularx}{\columnwidth}{@{}l Y@{}}
\toprule
\textbf{Parameter} & \textbf{Value} \\
\midrule
Supply & 3.3--5\,V (module) \\
Logic & 3.3\,V \\
LEDs & Red + IR \\
Internal ADC & Yes \\
\bottomrule
\end{tabularx}
\end{table}

\textbf{Selection rationale.} Widely used and documented; integrated optics and analog front end; compact and wearable-friendly.

\textbf{Integration notes.} Must sit flush against the skin; minimize ambient light leakage; reduce motion artifacts through mechanical design.

\subsection{ADS1115}
\label{sec:ads1115}

\textbf{Function.} High-resolution external ADC. \textbf{Purpose in PSYCON:} primary interface for GSR measurement.

\Cref{tab:ads1115-specs} lists the key specifications of the ADS1115.

\begin{table}[H]
\centering
\caption{Key specifications of the ADS1115.}
\label{tab:ads1115-specs}
\small
\begin{tabularx}{\columnwidth}{@{}l Y@{}}
\toprule
\textbf{Parameter} & \textbf{Value} \\
\midrule
Resolution & 16-bit \\
Channels & 4 \\
Interface & I\textsuperscript{2}C \\
Supply & 2.0--5.5\,V \\
PGA & Yes \\
\bottomrule
\end{tabularx}
\end{table}

\textbf{Default address.} \texttt{0x48}, configurable up to \texttt{0x49}, \texttt{0x4A}, \texttt{0x4B}.

\textbf{Why it was selected.} Better linearity and effective resolution than the ESP32 ADC; programmable gain supports low-level analog signals; well supported by existing libraries.

\begin{designnote}
The ADS1115 is intended for the GSR analog front end. Its additional channels remain available for future analog sensors.
\end{designnote}

\subsection{GSR Subsystem}
\label{sec:gsr-subsystem}

\textbf{Purpose.} Measure electrodermal activity (skin conductance).

\textbf{Current hardware.} ADS1115; electrodes.

\begin{requirementbox}
\textbf{Current status:} prototype-level implementation.
\end{requirementbox}

\begin{futureworkbox}
\textbf{Planned improvements:}
\begin{itemize}
    \item Constant-current excitation source.
    \item Patient-protection resistor(s).
    \item RC low-pass filtering.
    \item Improved analog conditioning.
    \item Stable reference.
    \item Shielded wiring if necessary.
\end{itemize}
\end{futureworkbox}

\begin{engnote}
The quality of GSR measurements depends primarily on the analog front end rather than the ADC alone. A documented and validated front end will be important for research-quality data.
\end{engnote}

\subsection{MPU6050}
\label{sec:mpu6050}

\textbf{Function.} Six-axis inertial measurement unit.

\textbf{Measurements.} 3-axis acceleration; 3-axis angular velocity.

\textbf{Interface.} I\textsuperscript{2}C. \textbf{Address:} default \texttt{0x68}, alternate \texttt{0x69}.

\textbf{Uses in PSYCON.} Motion detection; activity recognition; identifying motion artifacts in physiological signals.

\textbf{Why it was selected.} Mature and inexpensive; widely supported; adequate performance for wearable motion sensing.

\subsection{MCP9808}
\label{sec:mcp9808}

\textbf{Function.} Digital temperature sensor. \textbf{Measurement:} skin temperature.

\textbf{Interface.} I\textsuperscript{2}C. \textbf{Default address:} \texttt{0x18}, configurable through address pins.

\Cref{tab:mcp9808-specs} lists the key specifications of the MCP9808.

\begin{table}[H]
\centering
\caption{Key specifications of the MCP9808.}
\label{tab:mcp9808-specs}
\small
\begin{tabularx}{\columnwidth}{@{}l Y@{}}
\toprule
\textbf{Parameter} & \textbf{Value} \\
\midrule
Accuracy & $\pm$\SI{0.25}{\celsius} (typical) \\
Supply & 2.7--5.5\,V \\
Resolution & Up to \SI{0.0625}{\celsius} \\
\bottomrule
\end{tabularx}
\end{table}

\textbf{Selection rationale.} Chosen for higher accuracy and stability compared with lower-cost digital temperature sensors.

\subsection{INMP441}
\label{sec:inmp441}

\textbf{Function.} Digital MEMS microphone. \textbf{Interface:} I\textsuperscript{2}S. \textbf{Measurements:} digital audio for speech analysis.

\Cref{tab:inmp441-specs} lists the specifications of the INMP441.

\begin{table}[H]
\centering
\caption{Specifications of the INMP441 digital microphone.}
\label{tab:inmp441-specs}
\small
\begin{tabularx}{\columnwidth}{@{}l Y@{}}
\toprule
\textbf{Parameter} & \textbf{Value} \\
\midrule
Supply & 1.8--3.3\,V \\
Output & 24-bit I\textsuperscript{2}S \\
Directivity & Omnidirectional \\
\bottomrule
\end{tabularx}
\end{table}

\textbf{Selection rationale.} Native I\textsuperscript{2}S interface; low noise; good compatibility with the ESP32.

\textbf{Integration notes.} Keep away from switching power circuitry; ensure the acoustic port is unobstructed; verify gain, clipping margin, and DMA stability in firmware.

\subsection{TEMT6000}
\label{sec:temt6000}

\textbf{Function.} Ambient light sensor. \textbf{Output:} analog voltage proportional to light intensity. \textbf{Purpose:} provide environmental context (e.g., lighting conditions) for data interpretation.

\Cref{tab:temt6000-elec} lists the electrical characteristics of the TEMT6000.

\begin{table}[H]
\centering
\caption{Electrical characteristics of the TEMT6000.}
\label{tab:temt6000-elec}
\small
\begin{tabularx}{\columnwidth}{@{}l Y@{}}
\toprule
\textbf{Parameter} & \textbf{Value} \\
\midrule
Supply & 3.3--5\,V \\
Output & Analog \\
\bottomrule
\end{tabularx}
\end{table}

\begin{designnote}
For higher precision, the TEMT6000 output may be routed through the ADS1115 rather than the ESP32's internal ADC.
\end{designnote}

\subsection{TP4056 USB-C Charger (HW-373 V1.2.1)}
\label{sec:tp4056}

\textbf{Function.} Single-cell Li-ion/LiPo charging module with battery protection.

\textbf{Confirmed features.} USB-C input; battery protection; overcharge protection; over-discharge protection; over-current protection; short-circuit protection.

\textbf{Expected components.} TP4056 charger IC; DW01A protection IC; FS8205A dual MOSFET. IC markings should be verified from the physical board.

\begin{warningbox}
\textbf{Limitations:}
\begin{itemize}
    \item No true power-path / load-sharing.
    \item Not recommended to operate the wearable while charging.
    \item Reverse-polarity protection is not guaranteed.
\end{itemize}
\end{warningbox}

\subsection{LiPo Battery}
\label{sec:lipo-battery}

\textbf{Model.} HLCY 651735P. \Cref{tab:lipo-specs} lists its specifications.

\begin{table}[H]
\centering
\caption{Specifications of the HLCY 651735P LiPo battery.}
\label{tab:lipo-specs}
\small
\begin{tabularx}{\columnwidth}{@{}l Y@{}}
\toprule
\textbf{Parameter} & \textbf{Value} \\
\midrule
Chemistry & Lithium polymer \\
Nominal voltage & 3.7\,V \\
Fully charged & 4.2\,V \\
Capacity & 500\,mAh \\
\bottomrule
\end{tabularx}
\end{table}

\textbf{Configuration.} Wrist Module: 2 $\times$ 500\,mAh in parallel, total nominal capacity $\approx$1000\,mAh. Audio Module: 1 $\times$ 500\,mAh.

\begin{warningbox}
\textbf{Safety notes:}
\begin{itemize}
    \item Never connect the wrist batteries in series.
    \item Inspect for swelling, punctures, or damaged insulation before use.
    \item Disconnect charging before attaching GSR electrodes to a participant.
\end{itemize}
\end{warningbox}

\subsection{Hardware Compatibility Summary}
\label{sec:hardware-compatibility}

\Cref{tab:hardware-compatibility} summarizes pairwise compatibility between the ESP32 controller and each connected sensor, together with bus-level conflict checks.

\begin{table*}[t]
\centering
\caption{Hardware compatibility summary.}
\label{tab:hardware-compatibility}
\small
\begin{tabularx}{\textwidth}{@{}Y c@{}}
\toprule
\textbf{Component Pair / Check} & \textbf{Status} \\
\midrule
ESP32 $\leftrightarrow$ MAX30102 & \compatTag \\
ESP32 $\leftrightarrow$ MPU6050 & \compatTag \\
ESP32 $\leftrightarrow$ ADS1115 & \compatTag \\
ESP32 $\leftrightarrow$ MCP9808 & \compatTag \\
ESP32 $\leftrightarrow$ INMP441 & \compatTag \\
ESP32 $\leftrightarrow$ TEMT6000 & \compatTag \\
I\textsuperscript{2}C address conflicts & \noneTag \\
Logic-level conflicts & \noneTag\ (3.3\,V operation) \\
\bottomrule
\end{tabularx}
\end{table*}

\subsection{Remaining Hardware Validation}
\label{sec:remaining-hardware-validation}

\begin{requirementbox}
\textbf{The following still require experimental verification before final sign-off:}
\begin{itemize}
    \item ESP32 board revision and regulator identification.
    \item TP4056 PROG resistor value and actual charge current.
    \item Battery discharge characteristics under load.
    \item GSR analog front-end implementation and calibration.
    \item Current consumption in all operating modes.
    \item Thermal performance during charging and continuous operation.
    \item Brownout behavior and low-battery handling.
    \item Final wiring verification after assembly.
\end{itemize}
\end{requirementbox}

\subsection{Chapter Status}
\label{sec:ch4-status}

\begin{validationbox}
\textbf{Completed for this chapter:}
\begin{itemize}
    \item Full hardware inventory
    \item Detailed description of every module
    \item Electrical specifications
    \item Selection rationale
    \item Integration notes
    \item Power subsystem
    \item Safety considerations
    \item Compatibility review
    \item Remaining validation tasks
\end{itemize}
\end{validationbox}
